\documentclass[11pt]{article}
\usepackage[margin=1in]{geometry}
\usepackage{amsmath,amssymb,amsfonts,mathtools}
\usepackage{array,booktabs,enumitem,graphicx,xcolor,hyperref}
\usepackage[T1]{fontenc}
\usepackage[utf8]{inputenc}
\title{Supersymmetric Gelfand-Dickey Algebra}
\author{Katri Huitu and Dennis Nemeschansky}
\date{}
\begin{document}
\maketitle
\begin{abstract}

We study the classical version of supersymmetric \$W\$-algebras. Using the second Gelfand-Dickey Hamiltonian structure we work out in detail \$W\_2\$ and \$W\_3\$-algebras.

\end{abstract}

\section{Introduction}

.30in K. Huitu math expression Supported by a grant from the Finnish Cultural Foundation and D. Nemeschansky math expression Supported by the Department of Energy grant DE-FG03-84ER-40168 Department of Physics University of Southern California Los Angeles, CA 90089 .3in We study the classical version of supersymmetric math expression -algebras. Using the second Gelfand-Dickey Hamiltonian structure we work out in detail math expression and math expression -algebras.

\section{Method}

Above we have used math expression for multiplication to distinguish math expression from math expression . The coefficients math expression in eq. \textbackslash{} are the superbinomial coefficients, j k \textbackslash{} = \textbackslash{} \& math expression for math expression and for math expression mod math expression \& math expression for math expression , math expression mod math expression \textbackslash{} \textbackslash{} ,

\begin{equation}
\label{eq:compact}
L(\theta) = \sum_{i=1}^{n} \ell(x_i, \theta)
\end{equation}
Equation~\ref{eq:compact} is included to keep the sample paper-like while remaining small.

\section{Conclusion}

\& f(z, ) exp ( - 2 n \textasciicircum{}z u\_ n-1 (z\_1, ) dz\_1 ) f(z, ) \& f(z, ) exp (- 2 n-1 \textasciicircum{}z \textasciicircum{} z\textasciicircum{} D\_ z\_1, u\_ n-1 (z\_1, ) dz\_1 dz\textasciicircum{} ) f(z, ) \textbackslash{} . we can eliminate the coefficient math expression from the expansion of the differential operator . Therefore the differential operator math expression has the form L \textbackslash{} = \textbackslash{} D\textasciicircum{}n + u\_ n-2 (z, )D\textasciicircum{} n-2 + + u\_ 1 (z, ) D + u\_0 \textbackslash{} \textbackslash{} . Next let us consider the behavior of the differential operator under a superconformal transformation. A super analytic map math expression transforms the superderivative according to, e.g. \textbackslash{}
\end{document}
